% ==============================================================================
% ch05.tex -- Chapter 5: Regularity Theorems
%   Source: book pp. 129--151 (physical pp. 140--162)
% ==============================================================================
\chapter{Regularity theorems}
\label{ch:5}
\markboth{5.\ \MakeUppercase{Regularity theorems}}{}

% ------------------------------------------------------------------------------
% Section 5.1
% ------------------------------------------------------------------------------
\section{First Main Theorem}
\label{sec:5.1}

In this chapter, we study the regularity properties of the first three of the four basic moduli problems discussed in Chapter III (i.e., $\Gamma(N)$-structures, $\Gamma_1(N)$-structures, balanced $\Gamma_1(N)$-structures and $\Gamma_0(N)$-structures on variable elliptic curves $E/S$). For brevity, we will denote these moduli problems on $\Ell$ by the symbols
\[
[\Gamma(N)],\quad [\Gamma_1(N)],\quad [\mathrm{bal.}\ \Gamma_1(N)],\quad [\Gamma_0(N)].
\]

\begin{fmt}\label{thm:5.1.1}
Each of the four moduli problems $[\Gamma(N)]$, $[\Gamma_1(N)]$, $[\mathrm{bal.}\ \Gamma_1(N)]$, and $[\Gamma_0(N)]$ is relatively representable over $\Ell$. Each is finite and flat over $\Ell$ of constant rank $\ge 1$, and regular (necessarily of dimension two). Each tensored with $\Z[1/N]$ is finite etale over $\Ell/\Z[1/N]$.
\end{fmt}

% ------------------------------------------------------------------------------
% Section 5.2
% ------------------------------------------------------------------------------
\section{Axiomatics}
\label{sec:5.2}

By \ref{thm:3.7.1}, the theorem is true over $\Z[1/N]$. In view of the ``factorization into prime powers'' \ref{lem:3.5.1} of all the moduli problems involved, it suffices to treat the case when $N$ is a power of a single prime $p$.

To clarify the ideas which enter into the proof, we will explain the general mechanism (Deligne's ``homogeneity principle'') which underlies it (cf.\ \cite[Introduction]{DeRa}). Thus we consider any moduli problem $\mathcal{P}$ which satisfies the following axioms with respect to a fixed prime number $p$.

\smallskip
\noindent\textbf{Reg. 1.}\ \ $\mathcal{P}$ is relatively representable and finite over $\Ell$.

\smallskip
\noindent\textbf{Reg. 2.}\ \ $\mathcal{P}\otimes\Z[1/p]$ is finite etale over $\Ell/\Z[1/p]$.

\smallskip
\noindent\textbf{Reg. 3.}\ \ $\mathcal{P}$ depends only on the underlying $p$-divisible group, in the sense that for any two elliptic curves over a common base, say $E/S$ and $E'/S$, and any isomorphism of their $p$-divisible groups over $S$
\[
E[p^\infty] \xrightarrow{\ \sim\ } E'[p^\infty],
\]
there exists an isomorphism of $S$-schemes
\[
\mathcal{P}_{E'/S} \xrightarrow{\ \sim\ } \mathcal{P}_{E/S}
\]
(e.g., any $\mathcal{P}$ of $p$-power level in the sense of \ref{sec:4.14} satisfies Reg. 3).

\smallskip
\noindent\textbf{Reg. 4.}\ \ Let $k$ be an algebraically closed field of characteristic $p$, $E_0/k$ a supersingular elliptic curve, and $E/W(k)[[T]]$ its universal formal deformation. Then
\begin{enumerate}[label=(4\Alph*)]
\item $\mathcal{P}(E_0/k)$ consists of a single element, and
\item the (consequently local) finite $W[[T]]$-scheme
\[
\mathcal{P}_{E/W[[T]]}
\]
is the spectrum of a regular two-dimensional local ring.
\end{enumerate}

\begin{art}\label{thm:5.2.1}
Any moduli problem $\mathcal{P}$ satisfying the axioms Reg.~1 through Reg.~4 is finite and flat over $\Ell$ of constant rank $\ge 1$, and regular (necessarily of dimension two).
\end{art}

\begin{proof}
Let $S$ be any representable moduli problem which is etale over $\Ell$. We must show that $\mathfrak{M}(S, \mathcal{P})$ is a regular two-dimensional scheme, and that it is finite and flat over $\mathfrak{M}(S)$. Because $\mathfrak{M}(S)$ is (by \ref{cor:4.7.1}) a smooth curve over $\Z$, it is itself a regular two-dimensional scheme. By Reg.~2, $\mathfrak{M}(S, \mathcal{P})$ becomes finite etale over $\mathfrak{M}(S)$ as soon as we invert $p$. In particular, then, $\mathfrak{M}(S, \mathcal{P})\otimes\Z[1/p]$ is itself a regular two-dimensional scheme, finite and flat over $\mathfrak{M}(S)\otimes\Z[1/p]$.

Therefore it suffices to show that for some odd prime $\ell \ne p$, the scheme $\mathfrak{M}(S, \mathcal{P})\otimes\Z[1/\ell]$ is two-dimensional regular, and is finite and flat over $\mathfrak{M}(S)\otimes\Z[1/\ell]$. For this, we may replace $S$ by any representable moduli problem $S_1$ which is etale and surjective over $\Ell/\Z[1/\ell]$ (e.g., by $S_1$ the ``naive'' full level $\ell$ moduli problem, carried by $Y(\ell)$) simply because we have cartesian diagrams
\[
\begin{tikzcd}[row sep=large, column sep=large]
& \mathfrak{M}(S, S_1, \mathcal{P}) \arrow[dl] \arrow[dd] \arrow[dr] & \\
\mathfrak{M}(S, \mathcal{P})[1/\ell] \arrow[d] & & \mathfrak{M}(S_1, \mathcal{P}) \arrow[d] \\
\mathfrak{M}(S)[1/\ell] \arrow[r, "\alpha"] & \mathfrak{M}(S, S_1) \arrow[r, "\beta"] & \mathfrak{M}(S_1)
\end{tikzcd}
\]
in which $\alpha$ is etale and surjective, and in which $\beta$ is etale.

Now consider the finite morphism
\[
\begin{tikzcd}[row sep=small]
\mathfrak{M}(S_1, \mathcal{P}) \arrow[d, "\pi"] \\
\mathfrak{M}(S_1)
\end{tikzcd}
\]
We must show that
\[
\left\{
\begin{array}{@{}l@{}}
\mathfrak{M}(S_1, \mathcal{P}) \text{ is a regular scheme} \\
\mathfrak{M}(S_1, \mathcal{P}) \text{ is (finite and) flat over } \mathfrak{M}(S_1).
\end{array}
\right.
\]
(Recall that by Reg.~1, $\mathfrak{M}(S_1, \mathcal{P})$ is finite over $\mathfrak{M}(S_1)$.) Once these points are established, $\mathfrak{M}(S_1, \mathcal{P})$ will necessarily be a two-dimensional regular scheme, because $\mathfrak{M}(S_1)$ is itself regular two-dimensional, being a smooth curve over $\Z$.

Consider the open set $\mathfrak{U} \subset \mathfrak{M}(S_1)$ defined by
\[
\begin{array}{@{}l@{}}
\mathfrak{U} = \text{those points } y \text{ of } \mathfrak{M}(S_1) \text{ such that for all} \\
\quad\text{points } x \text{ of } \mathfrak{M}(S_1, \mathcal{P}) \text{ with } \pi(x) = y, \text{ the local ring} \\
\quad\cO_{\mathfrak{M}(S_1, \mathcal{P}), x} \text{ is regular, and is flat over the local ring} \\
\quad\cO_{\mathfrak{M}(S_1), y}.
\end{array}
\]

This is indeed an open set, for its complement is the union of the two closed sets in $\mathfrak{M}(S_1)$ which are the images by the finite (and hence proper) map $\pi$ of the closed subsets of $\mathfrak{M}(S_1, \mathcal{P})$ at which $\mathfrak{M}(S_1, \mathcal{P})$ is not regular, respectively is not flat over $\mathfrak{M}(S_1)$. We must show that $\mathfrak{U} = \mathfrak{M}(S_1)$.

Because $\mathfrak{M}(S_1)$ is of finite type over $\Z$, it suffices to show that $\mathfrak{U}$ contains every closed point of $\mathfrak{M}(S_1)$. By Reg.~2, we know that $\mathfrak{U}$ contains all of $\mathfrak{M}(S_1)\otimes\Z[1/p]$, so it suffices to show that $\mathfrak{U}$ contains every closed point of $\mathfrak{M}(S_1)$ of residue characteristic $p$, i.e., every closed point of $\mathfrak{M}(S_1)\otimes\F_p$.

Recall that a closed point of $\mathfrak{M}(S_1)\otimes\F_p$ is said to be ordinary (respectively supersingular) if the corresponding elliptic curve with $S_1$-structure over the residue field is an ordinary (respectively supersingular) elliptic curve. The set of supersingular points of $\mathfrak{M}(S_1)\otimes\F_p$ is finite and non-empty, and every open neighborhood of a supersingular point contains ordinary points (\ref{thm:2.9.4}, \ref{cor:4.7.2}).

We will show that the open set $\mathfrak{U}$ has the following ``homogeneity properties.''
\begin{enumerate}[label=(H\arabic*)]
\item if $\mathfrak{U}$ contains one supersingular point of $\mathfrak{M}(S_1)\otimes\F_p$, it contains all supersingular points of $\mathfrak{M}(S_1)\otimes\F_p$.
\item if $\mathfrak{U}$ contains one ordinary point of $\mathfrak{M}(S_1)\otimes\F_p$, it contains all ordinary points of $\mathfrak{M}(S_1)\otimes\F_p$.
\end{enumerate}
If we grant that $\mathfrak{U}$ satisfies (H1) and (H2), then we are reduced to showing that $\mathfrak{U}$ contains some supersingular point of $\mathfrak{M}(S_1)\otimes\F_p$ [by (H1), $\mathfrak{U}$ will contain all the supersingular points; being open, it also contains some ordinary point, so by (H2) it contains all ordinary points].

We now explain why (H1) and (H2) are true. By definition, a closed point $y$ of $\mathfrak{M}(S_1)$ lies in $\mathfrak{U}$ if and only if for each of the finitely many closed points $x$ of $\mathfrak{M}(S_1, \mathcal{P})$ lying over it, we have:
\[
\text{the local ring } \cO_{\mathfrak{M}(S_1,\mathcal{P}),x} \text{ is regular, and it is flat over } \cO_{\mathfrak{M}(S_1),y}.
\]
We first reduce to the case of an algebraically closed residue field, by passing to the completion of the strict henselization. Let $k$ be an algebraic closure of $\F_p$, and $W = W(k)$ its ring of Witt vectors. For any scheme $X$ of finite type over $\Z$, there are natural bijections
\[
\begin{array}{@{}l@{}}
\left\{\text{closed points } x \text{ of } X\otimes W \text{ with residue field } k\right\} \longleftrightarrow \\[6pt]
\left\{k\text{-valued points } x \text{ of } X\right\} \longleftrightarrow \\[6pt]
\left\{
\begin{array}{@{}l@{}}
\text{pairs } (x_0, i) \text{ where } x_0 \text{ is a closed point of } X \text{ with residue} \\
\text{characteristic } p \text{ and } i \text{ is an inclusion of the finite residue} \\
\text{field } \F_p(x_0) \text{ into } k
\end{array}
\right\}.
\end{array}
\]
Under this bijection the complete local rings are related by
\[
\widehat{\cO}_{X\otimes W, x} \ \simeq \ \widehat{\cO^{h.s.}_{X, X_p}}.
\]
Recall that for any noetherian local ring $A$, both its strict henselization $A^{h.s.}$ and its completion $\hat{A}$ are noetherian local rings which are faithfully flat over $A$, and which are regular if and only if $A$ is regular (for $\hat{A}$ this is standard, for $A^{h.s.}$ it's \cite[IV, 18.8.8, 18.8.13]{EGA}). Applying this to the local rings at closed points of both $\mathfrak{M}(S_1)$ and $\mathfrak{M}(S_1, \mathcal{P})$, it follows that a closed point $y_p$ of $\mathfrak{M}(S_1)$ of residue characteristic $p$ lies in $\mathfrak{U}$ if and only if for any $k$-valued point $y$ of $\mathfrak{M}(S_1)\otimes W$ lying over it, we have the following condition:
\[
\left\{
\begin{array}{@{}l@{}}
\text{for all } k\text{-valued points } x \text{ of } \mathfrak{M}(S_1, \mathcal{P})\otimes W \text{ lying over } y, \text{ the} \\
\text{complete local ring } \widehat{\cO}_{\mathfrak{M}(S_1,\mathcal{P})\otimes W, x} \text{ is regular,} \\
\text{and it is flat over } \widehat{\cO}_{\mathfrak{M}(S_1)\otimes W, y}.
\end{array}
\right.
\]
Because $\mathfrak{M}(S_1, \mathcal{P})$ is finite over $\mathfrak{M}(S_1)$, the spectrum of the direct sum (over all $x \to y$) of these complete local rings is just the scheme
\[
(*)\qquad \left(\mathfrak{M}(S_1, \mathcal{P})\otimes W\right) \ \times_{\mathfrak{M}(S_1)\otimes W} \ \Spec\left(\widehat{\cO}_{\mathfrak{M}(S_1)\otimes W, y}\right).
\]
Consider the elliptic curve $E_0/k$ underlying $y$ viewed as a $k$-valued point of $\mathfrak{M}(S_1)$, and let $E/W[[T]]$ be its universal formal deformation.

Because $\mathfrak{M}(S_1)$ is etale over $\Ell$, we have an isomorphism of complete local rings
\[
W[[T]] \ \simeq \ \widehat{\cO}_{\mathfrak{M}(S_1)\otimes W, y},
\]
with $E/W[[T]]$ identified with the inverse image of the universal curve over $\mathfrak{M}(S_1)$. Therefore the scheme $(*)$ in question is none other than
\[
\mathcal{P}_{E/W[[T]]}.
\]
The condition that $y$ lie in $\mathfrak{U}$ is thus a condition on the $W[[T]]$-scheme $\mathcal{P}_{E/W[[T]]}$ (i.e., regular scheme, flat over $W[[T]]$).

By the Serre-Tate theorem, the $p$-divisible group of $E/W[[T]]$ is the universal formal deformation of the $p$-divisible group of $E_0/k$. By Reg.~3, the $W[[T]]$-scheme $\mathcal{P}_{E/W[[T]]}$ depends only on the $p$-divisible group of $E/W[[T]]$. Therefore by Serre-Tate the $W[[T]]$-scheme $\mathcal{P}_{E/W[[T]]}$ depends only on the $p$-divisible group of $E_0/k$. Because $k$ is algebraically closed, the $p$-divisible group of $E_0/k$ is necessarily $k$-isomorphic to either $\Q_p/\Z_p \times \mu_{p^\infty}$, corresponding to $E_0/k$ ordinary, or to the unique one-parameter formal Lie group over $k$ of height two, corresponding to $E_0/k$ supersingular.

Therefore up to $W[[T]]$-isomorphism, there are only two distinct schemes $\mathcal{P}_{E/W[[T]]}$, one corresponding to $E_0/k$ ordinary, the other to $E_0/k$ supersingular.

Therefore the open set $\mathfrak{U}$ contains either all ordinary (respectively all supersingular) points in characteristic $p$, or it contains none; the conditions (H1) and (H2) are satisfied.

To complete the proof, it remains to establish that $\mathfrak{U}$ contains some supersingular point of characteristic $p$. But by Reg.~4, the open set $\mathfrak{U}$ contains all supersingular points of characteristic $p$.

Finally, we recall \ref{lem:4.12.1} that once we know $\mathcal{P}$ is finite and flat over $\Ell$, its rank is necessarily constant.

This concludes the proof of the Axiomatic Regularity Theorem.
\end{proof}

In practice, it is convenient to replace axiom Reg.~4B by a more computationally accessible condition Reg.~4B bis:

\begin{customprop}{5.2.2}\label{prop:5.2.2}
Let $\mathcal{P}$ be a moduli problem satisfying Reg.~1 through Reg.~3 and Reg.~4A. Then $\mathcal{P}$ satisfies Reg.~4B if and only if it satisfies
\[
\begin{array}{@{}l@{\quad}l@{}}
\textbf{\itshape Reg.\ 4B bis} & \mathcal{P}_{E/W[[T]]} \text{ is the spectrum of a local ring whose} \\
& \text{maximal ideal is generated by two elements.}
\end{array}
\]
\end{customprop}

\begin{proof}
Clearly 4B implies 4B bis. We must establish the converse. Write $\mathcal{P}_{E/W[[T]]} = \Spec(A)$. Then $A$ is local by 4A, and is finite over $W[[T]]$ by Reg.~1. By Reg.~2, $A\otimes\Z[1/p]$ is finite etale over $(W[[T]])\otimes\Z[1/p]$. Therefore the finite morphism
\[
\Spec(A) \longrightarrow \Spec(W[[T]])
\]
is surjective, because its image is closed (the map being finite, hence proper) and its image contains the open dense set $\Spec((W[[T]])\otimes\Z[1/p])$. Therefore the local ring $A$ must have dimension at least two. So if its maximal ideal is generated by two elements, $A$ must be a regular local ring of dimension two. It is then automatically flat over $W[[T]]$, because any finite morphism between regular schemes of the same dimension is automatically flat \cite[V, 3.6]{AK1}.
\end{proof}

% ------------------------------------------------------------------------------
% Section 5.3
% ------------------------------------------------------------------------------
\section{End of the proof}
\label{sec:5.3}

In this section, we will complete the proof of the First Main Theorem for the three moduli problems
\[
[\Gamma(N)],\quad [\Gamma_1(N)],\quad\text{and}\quad [\mathrm{bal.}\ \Gamma_1(N)].
\]
As already noted, it suffices to treat the case $N =$ a power of a prime $p$, say $N = p^n$ with $n \ge 1$.

Each of these problems, with $N = p^n$, satisfies Reg.~1 through Reg.~3 (\ref{sec:3.6}, \ref{thm:3.7.1}, \ref{item:4.14.3}). We must show that each also satisfies Reg.~4A and Reg.~4B bis. We begin with 4A.

\begin{klem}\label{lem:5.3.1}
Each of the problems
\[
\mathcal{P} = [\Gamma(p^n)],\quad [\Gamma_1(p^n)],\quad [\mathrm{bal.}\ \Gamma_1(p^n)]
\]
has $\mathcal{P}(E_0/k) = $ a single element for any supersingular elliptic curve over any field $k$ of characteristic $p$.
\end{klem}

\begin{proof}
Because $E_0$ is supersingular, the group-scheme $E_0[p^n]$ over $k$ is local for every $n$, and the unique closed subgroup-scheme of $E_0$ of any given rank $p^r$ is the kernel of the $r$-fold Frobenius morphism
\[
F^r : E_0 \longrightarrow E_0^{(p^r)}.
\]
In particular, we have
\[
\Ker(F^{2n}) = \Ker(p^n)
\]
by uniqueness.

Consider first $\mathcal{P} = [\Gamma(p^n)]$. Then $\mathcal{P}(E_0/k)$ is the set of pairs $P, Q$ of $k$-points of $E_0[p^n]$ which form a Drinfeld basis. Because $E_0[p^n]$ is local, and $k$ is a field, we have $P = Q = 0$. This is in fact a Drinfeld basis of $E_0[p^n] = \Ker(p^n) = \Ker(F^{2n})$, simply because for $P = Q = 0$ we have the equality of Cartier divisors
\[
\sum_{a, b \bmod p^n} [aP + bQ] = p^{2n}[0] = \Ker(F^{2n}).
\]

Consider next $\mathcal{P} = [\Gamma_1(p^n)]$. Then $\mathcal{P}(E_0/k)$ is the set of $k$-valued points $P$ of $E_0[p^n]$ of ``exact order $p^n$.'' Again the only possibility is $P = 0$, which does in fact have ``exact order $p^n$,'' because it generates $\Ker(F^n)$.

Finally for $\mathcal{P} = [\mathrm{bal.}\ \Gamma_1(p^n)]$, the set $\mathcal{P}(E_0/k)$ is the set of triples $(K, P, P')$ with $K$ a finite flat $k$-subgroup-scheme of $E_0[p^n]$ of rank $p^n$, $P \in K(k)$ a generator of $K$, and $P'$ a $k$-valued generator of the quotient group $K' = E_0[p^n]$ mod $K$. Then $K$ and $K'$ are local, so $P = 0$ and $Q = 0$, and $K$ must be $\Ker(F^n)$.
\end{proof}

We now turn to the condition 4B bis.

\begin{kthm}\label{thm:5.3.2}
Let $k$ be an algebraically closed field of characteristic $p$, $E_0/k$ a supersingular elliptic curve, $E/W[[T]]$ its universal formal deformation. Then for each of the three moduli problems
\[
\mathcal{P} = [\Gamma(p^n)],\quad [\Gamma_1(p^n)],\quad [\mathrm{bal.}\ \Gamma_1(p^n)],
\]
we have
\[
\mathcal{P}_{E/W[[T]]} = \Spec(A),
\]
with $A$ a local ring whose maximal ideal is generated by two elements.
\end{kthm}

\begin{proof}
By 4A, we know that $A$ is a local ring, finite over $W[[T]]$, hence complete. Consider the elliptic curve $\mathcal{E}/A$ obtained from $E/W[[T]]$ by the extension of scalars $W[[T]] \to A$. Let $\alpha$ be the tautological level $\mathcal{P}$ structure on $\mathcal{E}/A$. Then $(\mathcal{E}/A, \alpha)$ pro-represents the functor on artin local $W(k)$ algebras with residue field $k$
\[
R \longmapsto
\begin{array}{@{}l@{}}
\text{triples } (E/R, \alpha, i) \text{ where } E/R \text{ is an elliptic curve,} \\
\alpha \in \mathcal{P}(E/R) \text{ is a level } \mathcal{P} \text{ structure, and } i \text{ is an} \\
\text{isomorphism of elliptic curves over } k, \\
E\otimes_R (R/\max(R)) \xrightarrow{\ \sim\ } E_0,
\end{array}
\]
which (necessarily) carries $\alpha$ to the unique element of $\mathcal{P}(E_0/k)$.

Suppose we are given two functions $f, g$ in $\max(A)$. Then in order that $f$ and $g$ generate $\max(A)$, it is necessary and sufficient that the following moduli-theoretic rigidity condition hold:

\smallskip
\noindent If $R$ is an artin local $W(k)$-algebra,\footnote{With residue field $k$.} and if $\varphi : A \to R$ is a homomorphism with $\varphi(f) = \varphi(g) = 0$, then the induced triple $(E/R, \alpha, i)$ is constant in the sense that $R$ is a $k$-algebra and this triple comes from the constant triple $(E_0/k, \alpha_0, \mathrm{id})$ by extension of scalars $k \to R$.
\smallskip

We remark that in this condition, it is sufficient that $(E/R, \alpha)$ be constant, for by rigidity the only hom's between constant abelian schemes are themselves constant.

\medskip
\noindent\emph{Case I.}\ $\mathcal{P} = [\Gamma(p^n)]$. Choose a parameter $X$ for the formal group of $\mathcal{E}/A$. The universal $\mathcal{P}$-structure on $\mathcal{E}/A$ is a pair of points $P, Q$ in $\mathcal{E}(A)$ which form a Drinfeld $p^n$-basis. Because $A$ is complete local, and $E_0/k$ is supersingular, the points $P, Q$ both lie in the formal group. Therefore we may speak of the $X$-coordinates of these points,
\[
X(P),\ X(Q) \in \max(A).
\]
We claim that the two functions $f = X(P)$, $g = X(Q)$, generate $\max(A)$. This amounts to the following rigidity assertion:

\begin{pitemnum}{5.3.2.1}\label{lem:5.3.2.1}
(Rigid I). Let $k$ be an algebraically closed field of characteristic $p$. If $R$ is an artin local $W(k)$-algebra, and if $E/R$ admits $(0,0)$ as a Drinfeld $p^n$-basis, then $R$ is a $k$-algebra and $E/R$ is constant.
\end{pitemnum}

\medskip
\noindent\emph{Case II.}\ $\mathcal{P} = [\Gamma_1(p^n)]$. Again choose a parameter $X$ for the formal group of $\mathcal{E}/A$. The universal $\mathcal{P}$-structure on $\mathcal{E}/A$ is a point $P$ in $\mathcal{E}(A)$ of ``exact order $p^n$,'' and just as above $P$ lies in the formal group, so we may speak of its $X$-coordinate
\[
X(P) \in \max(A).
\]
We also have ``the'' parameter $T$ of $W[[T]]$ (strictly speaking, any element of $W[[T]]$ which, together with $p$, generates the maximal ideal of $W[[T]]$). By means of $W[[T]] \to A$, we may speak of $T$ as lying in $\max(A)$. We claim that the two functions
\[
X(P),\ T \in \max(A)
\]
generate $\max(A)$. This amounts to the following rigidity assertion:

\begin{pitemnum}{5.3.2.2}\label{lem:5.3.2.2}
(Rigid II). Let $k$ be an algebraically closed field of characteristic $p$. If $R$ is an artin local $W(k)$-algebra,\footnote{With residue field $k$.} and if $E/R$ admits $0$ as a point of ``exact order $p^n$,'' then $R$ is a $k$-algebra, i.e., $p = 0$ in $R$.
\end{pitemnum}

For if we grant that $p = 0$ in $R$, then the vanishing of the $T$-invariant of $E/R$ means precisely that $E/R$ is constant.

\medskip
\noindent\emph{Case III.}\ $\mathcal{P} = [\mathrm{bal.}\ \Gamma_1(p^n)]$. This time we choose a parameter $X$ for the formal group of $\mathcal{E}/A$, and we choose a parameter $X'$ for the formal group of $\mathcal{E}'/A$, where $\mathcal{E}'$ is the quotient of $\mathcal{E}$ by the cyclic subgroup-scheme $K$ of rank $p^n$ of $\mathcal{E}$ which is part of the data of the $\mathcal{P}$-structure. All in all, we have
\[
P \in \mathcal{E}(A),\quad P' \in \mathcal{E}'(A),
\]
where $P$ has ``exact order $p^n$'' on $\mathcal{E}/A$ and generates $\Ker \pi$, and where $P'$ has ``exact order $p^n$'' on $\mathcal{E}'/A$ and generates $\Ker(\pi^t)$. As in Cases I and II, the points $P$ and $P'$ lie in the formal groups of $\mathcal{E}$ and $\mathcal{E}'$ respectively, so we can speak of their $X$ and $X'$-coordinates:
\[
X(P),\ X'(P') \in \max(A).
\]
We claim that these two functions generate $\max(A)$. This amounts to the following rigidity assertion:

\begin{pitemnum}{5.3.2.3}\label{lem:5.3.2.3}
(Rigidity III). Let $k$ be an algebraically closed field of characteristic $p$. If $R$ is an artin local $W(k)$-algebra and if $E$ and $E'$ are elliptic curves over $R$ joined by dual $p^n$-isogenies
\[
\begin{tikzcd}[row sep=small, column sep=large]
E \arrow[r, shift left, "\pi"] & E' \arrow[l, shift left, "\pi^t"]
\end{tikzcd}
\]
whose kernels are both cyclic and generated by their zero-sections, then $R$ is a $k$-algebra, $E/R$ is constant, and $\Ker(\pi) = \Ker(F^n)$.
\end{pitemnum}

We now turn to proving the rigidity assertions I, II, III. For ease of later reference, we will give an analysis of a slightly more general situation, which is useful in other contexts as well (compare our discussion of roots of unity \ref{sec:1.12}).

\begin{kprop}\label{prop:5.3.3}
Let $R$ be an arbitrary ring, $C/R$ a smooth commutative one-dimensional group-scheme over $R$ (cf.\ \ref{item:1.4.1}), $p$ a prime number, $n \ge 1$ an integer. Suppose that the zero section $0 \in C(R)$ is a point of ``exact order $p^n$'' in $C/R$, i.e., that the effective Cartier divisor $p^n[0]$ in $C/R$ is a subgroup-scheme. Then $p = 0$ in $R$, and the subgroup-scheme $p^n[0]$ is none other than $\Ker(F^n)$.
\end{kprop}

\medskip
\noindent\emph{Proof.}\ The zero-section and the Cartier divisor $p^n[0]$ both lie in the formal group of $C/R$. Zariski locally on $R$, we may choose a parameter $X$ for the formal group. Once this is done, our proposition results from the following one.

\begin{kprop}\label{prop:5.3.4}
Let $R$ be an arbitrary ring, and
\[
G(X, Y) = X + Y + \cdots
\]
a one-parameter commutative formal group law over $R$. Suppose that for some prime power $p^n$, the equation
\[
X^{p^n} = 0
\]
defines a subgroup-scheme of $G$. Then $p = 0$ in $R$, and this subgroup-scheme is equal to $\Ker(F^n)$.
\end{kprop}

\begin{proof}
Once we prove $p = 0$ in $R$, the Frobenius morphism is defined, and we visibly have that $\Ker(F^n)$ is defined by $X^{p^n} = 0$. To prove that $p = 0$ in $R$, we argue as follows. The hypothesis means that for any $R$-algebra $B$, and any elements $x, y \in B$, we have the implication
\[
x^{p^n} = y^{p^n} = 0 \text{ in } B \implies (G(x,y))^{p^n} = 0 \text{ in } B.
\]
Considering the universal situation
\[
B = R[[X,Y]]/(X^{p^n}, Y^{p^n}),
\]
we see that inside $R[[X,Y]]$, we have
\[
(G(X,Y))^{p^n} \in (X^{p^n}, Y^{p^n}),
\]
say
\[
G(X,Y)^{p^n} = X^{p^n}\cdot A(X,Y) + Y^{p^n}\cdot B(X,Y).
\]
Comparing terms of total degree $p^n$, we obtain
\[
(X+Y)^{p^n} = X^{p^n}\cdot A(0,0) + Y^{p^n}\cdot B(0,0).
\]
Comparing coefficients of $X^{p^n}$ and $Y^{p^n}$, we see that $A(0,0) = B(0,0) = 1$, whence
\[
(X+Y)^{p^n} = X^{p^n} + Y^{p^n} \text{ in } R[[X,Y]].
\]
Explicitly, this means that all the intermediate binomial coefficients vanish in $R$:
\[
\binom{p^n}{i} = 0 \text{ in } R \quad\text{for } i = 1, \ldots, p^n - 1.
\]
Taking $i = 1$, we see that $p^n = 0$ in $R$. In particular, $R$ is a $\Z_p$-algebra. Taking $i = p^{n-1}$, and using the standard fact that
\[
\binom{p^n}{p^{n-1}} = p \times (\text{an integer prime to } p),
\]
we see that $p = 0$ in $R$.
\end{proof}

\begin{pitem}\label{lem:5.3.5}
It is now a simple matter to deduce the rigidity assertions I, II, III, (i.e., \ref{lem:5.3.2.1}-2-3). Assertion II is just the proposition itself, applied to $E/R$. In Assertion III, the proposition applied directly to $E/R$ shows that $p = 0$ in $R$ and that $\Ker(\pi) = \Ker(F^n)$. To finish the proof of III, we must show that $E/R$ is constant. This itself results from Assertion I, for it follows from \ref{prop:1.11.3}, applied to the commutative diagram
\begin{equation}\label{eq:5.3.5.1}
\begin{tikzcd}[row sep=large, column sep=small]
0 \arrow[r] & K \arrow[r] & E[p^n] \arrow[r] & K' \arrow[r] & 0 \\
0 \arrow[r] & \Z/p^n\Z \arrow[r] \arrow[u, "0"'] & (\Z/p^n\Z)^2 \arrow[r] \arrow[u, "0"'] & (\Z/p^n\Z) \arrow[r] \arrow[u, "0"'] & 0
\end{tikzcd}
\end{equation}
that $(0,0)$ is a Drinfeld $p^n$-basis of $E/R$.
\end{pitem}

\begin{pitem}\label{lem:5.3.6}
To prove Assertion I, we argue as follows. If $(0,0)$ is a Drinfeld $p^n$-basis of $E/R$, then we have an equality of Cartier divisors
\[
p^{2n}[0] = \Ker(p^n).
\]
In particular, $p^{2n}[0]$ is a subgroup, so Proposition \ref{prop:5.3.3} (with $E/R$, $p^{2n}$) shows that $p = 0$ in $R$, and that
\[
\Ker(p^n) = \Ker(F^{2n}).
\]
To show that $E/R$ is constant, we argue as follows. We have exact sequences of $R$-group-schemes
\[
\begin{tikzcd}[row sep=small, column sep=small]
0 \arrow[r] & \Ker(p^n) \arrow[r] & E \arrow[r, "p^n"] & E \arrow[r] & 0 \\
0 \arrow[r] & \Ker(F^{2n}) \arrow[r] & E \arrow[r, "F^{2n}"] & E^{(p^{2n})} \arrow[r] & 0.
\end{tikzcd}
\]
Because the kernels are equal, we have an $R$-isomorphism
\[
E \xrightarrow{\ \sim\ } E^{(p^{2n})}.
\]
Taking the Frobenius pull-backs of these isomorphisms, we obtain isomorphisms
\[
E \xrightarrow{\ \sim\ } E^{(p^{2n})} \xrightarrow{\ \sim\ } E^{(p^{4n})} \xrightarrow{\ \sim\ } E^{(p^{6n})} \ldots .
\]
But for all sufficiently large $r$, the elliptic curve
\[
E^{(p^r)}
\]
is constant, simply because $R$ is an artin local $k$-algebra, so that for all $r \gg 0$ the $p^r$-th power map on $R$ factors
\[
R \xrightarrow{\ F^r\ } k \longrightarrow R,
\]
(i.e., we recover $k$ as the subring of $p^r$-th powers for all $r \gg 0$).

This concludes the proof of the rigidity assertions I, II, III, and with them the proof of the First Main Theorem \ref{thm:5.1.1} for $[\Gamma(N)]$, $[\Gamma_1(N)]$, $[\mathrm{bal.}\ \Gamma_1(N)]$.
\end{pitem}
\end{proof}

The proof of the First Main Theorem \ref{thm:5.1.1} for $\Gamma_0(N)$ will be given in Chapter 6.

% ------------------------------------------------------------------------------
% Section 5.4
% ------------------------------------------------------------------------------
\section{Summary of parameters at supersingular points}
\label{sec:5.4}

$E_0/k$ is a supersingular elliptic curve over a perfect field $k$ of characteristic $p$, $E/W(k)[[T]]$ is its universal formal deformation, $X$ is a parameter for its formal group, $A$ is the complete local ring such that $\mathcal{P}_{E/W[[T]]} = \Spec(A)$. If $\mathcal{P}$ is bal-$\Gamma_1$ or $\Gamma_1$, then $E'/A$ is the quotient of $E\otimes A$ by the universal cyclic subgroup it carries, and $X'$ is a parameter for the formal group of $E'$ over $A$.

\begin{center}
\begin{tabular}{@{}c|c|c@{}}
\hline
Moduli problem & basic data & parameters for $A$ \\
\hline
$[\Gamma(p^n)]$ & Drinfeld $p^n$-basis $(P, Q)$ & $X(P),\ X(Q)$ \\
\hline
$[\mathrm{bal.}\ \Gamma_1(p^n)]$ & $\begin{array}{@{}l@{}} P\ ;\ E \overset{\pi}{\underset{\pi^t}{\rightleftarrows}} E'\ ;\ P' \\[2pt] P \text{ generates } \Ker \pi, \\ P' \text{ generates } \Ker \pi^t \end{array}$ & $X(P),\ X'(P')$ \\
\hline
$[\Gamma_1(p^n)]$ & point $P$ of exact order $p^n$ & $T,\ X(P)$. \\
\hline
\end{tabular}
\end{center}

% ------------------------------------------------------------------------------
% Section 5.5
% ------------------------------------------------------------------------------
\section{First applications}
\label{sec:5.5}

In this section, we will deduce some immediate consequences of the First Main Theorem for $\Gamma$, $\Gamma_1$, and balanced $\Gamma_1$.

\begin{kthm}\label{thm:5.5.1}
Let $N > 1$ be an integer, and let $S$ be a representable moduli problem which is etale over $\Ell$. Then for
\[
\mathcal{P} = [\Gamma(N)],\quad [\Gamma_1(N)],\quad [\mathrm{bal.}\ \Gamma_1(N)],
\]
we have
\begin{enumerate}[(1)]
\item $\mathfrak{M}(S, \mathcal{P})$ is a regular two-dimensional scheme, finite and flat over $\mathfrak{M}(S)$.
\item $\mathfrak{M}(S, \mathcal{P})\otimes\Z[1/N]$ is finite etale over $\mathfrak{M}(S)\otimes\Z[1/N]$.
\item $\mathfrak{M}(S, \mathcal{P})$ is flat over $\Z$.
\item $\mathfrak{M}(S, \mathcal{P})$ is the normalization of $\mathfrak{M}(S)$ in $\mathfrak{M}(S, \mathcal{P})\otimes\Z[1/N]$.
\end{enumerate}
\end{kthm}

\begin{proof}
Assertions (1) and (2) are simply ``mise pour memoire.'' Assertion (3) results from (1) because $\mathfrak{M}(S)$ is itself flat over $\Z$, being a smooth curve over $\Z$. Assertion (4) results from (1) and (2), because the normalization in question is the unique normal scheme finite over $\mathfrak{M}(S)$ which agrees with $\mathfrak{M}(S, \mathcal{P})$ outside of $N$.\footnote{i.e., over the open set where $N$ is invertible.}
\end{proof}

\begin{kthm}\label{thm:5.5.2}
Let $E/S$ be an elliptic curve, $N > 1$ an integer, and $P, Q \in E[N](S)$ a Drinfeld $N$-basis of $E/S$. Then
\begin{enumerate}[(1)]
\item the point $P$ has ``exact order $N$'' on $E/S$, i.e., $P$ is a $\Gamma_1(N)$-structure on $E/S$.
\item Let $K \subset E[N]$ denote the cyclic subgroup-scheme of rank $N$ generated by $P$. The image $Q'$ of $Q$ in the quotient elliptic curve $E'$ of $E$ modulo $K$ has ``exact order $N$'' on $E'$, and $Q'$ generates $K'$, the quotient of $E[N]$ by $K$. In other words, the data $(P, K, Q')$ is a balanced $\Gamma_1(N)$-structure on $E/S$.
\end{enumerate}
\end{kthm}

\begin{proof}
The theorem is obvious if $N$ is invertible on $S$, (cf.\ \ref{cor:3.7.2}). We will reduce to that case.

The question is f.p.p.f. local on $S$, so we may assume that some odd prime $\ell$ is invertible on $S$ and that $E/S$ admits a naive full level $\ell$ structure, which we abbreviate ``$S_1$-structure.'' By reduction to the universal case, we may assume that $S = \mathfrak{M}(S, \Gamma(N))$ and that $E/S$ is the corresponding universal curve, $(P, Q)$ its universal $\Gamma(N)$-structure. In particular, we may assume that the base $S$ is flat over $\Z$, and affine, say $S = \Spec(A)$.

Let us prove (1). We know that the $\Gamma_1(N)$-moduli problem $\mathcal{P}$ is relatively representable by a closed subscheme $\mathcal{P}_{E/A}$ of the affine $A$-scheme $\Hom(\Z/N\Z, E) = E[N]$:
\[
\mathcal{P}_{E/A} \hookrightarrow E[N].
\]
Pick a collection of functions $f_i$ in the affine ring of $E[N]$ which define the closed subscheme $\mathcal{P}_{E/A}$. We must show that, for all $i$,
\[
f_i(P) = 0 \text{ in } A.
\]
But the theorem is true as soon as we invert $N$, i.e., we know that
\[
f_i(P) = 0 \text{ in } A[1/N].
\]
Because $A$ is flat over $\Z$, we have $A \subset A[1/N]$, whence $f_i(P) = 0$ in $A$, as required.

Let us now prove (2). Let $K$ denote the cyclic subgroup-scheme of $E[N]$ ``generated'' by $P$, let $K'$ denote the quotient of $E[N]$ by $K$, and let $Q' \in K'(A)$ denote the image of the point $Q$. We must show that $Q'$ ``generates'' $K'$. But the scheme of generators of $K'$ is a closed subscheme of the affine $A$-scheme $\Hom(\Z/N\Z, K') = K'$:
\[
(\Z/N\Z)\text{-Gen}(K'/A) \hookrightarrow K'.
\]
Again, choosing a collection of functions $f_i$ in the affine ring of $K'$ which defines the closed subscheme of generators, we must show that
\[
f_i(Q') = 0 \text{ in } A.
\]
Because the theorem is true as soon as we invert $N$, we know that
\[
f_i(Q') = 0 \text{ in } A[1/N],
\]
so again by the flatness of $A$ over $\Z$, we conclude that $f_i(Q') = 0$ in $A$.
\end{proof}

\begin{kcor}\label{cor:5.5.3}
The constructions of the preceding theorem define natural morphisms of moduli problems
\[
[\Gamma(N)] \longrightarrow [\mathrm{bal.}\ \Gamma_1(N)] \longrightarrow [\Gamma_1(N)]
\]
which are finite and flat of the indicated degrees
\[
[\Gamma(N)] \xrightarrow{\ \text{degree } N\ } [\mathrm{bal.}\ \Gamma_1(N)] \xrightarrow{\ \text{degree } \varphi(N)\ } [\Gamma_1(N)].
\]
Concretely, for any representable moduli problem $S$, we have a natural diagram of morphisms, each of which is finite flat of the indicated degree
\[
\text{degree} = \#GL(2, \Z/N\Z)\;
\left\{\;
\begin{tikzcd}[row sep=large, column sep=small]
\mathfrak{M}(S, \Gamma(N)) \arrow[d, "\text{degree } N"] \\
\mathfrak{M}(S, \mathrm{bal.}\ \Gamma_1(N)) \arrow[d, "\text{degree } \varphi(N)"] \\
\mathfrak{M}(S, \Gamma_1(N)) \arrow[d, "{\text{degree } N^2\cdot\prod_{p\mid N}\left(1-\frac{1}{p^2}\right)}"] \\
\mathfrak{M}(S)
\end{tikzcd}
\right.
\]
\end{kcor}

\begin{proof}
That these morphisms are defined is the content of the preceding theorem. To see that they are finite and flat of the asserted degrees, we may reduce to the case when $S$ is etale over $\Ell$, and $\mathfrak{M}(S)$ is connected. In this case all the schemes involved are regular two-dimensional schemes, finite over $\mathfrak{M}(S)$. Therefore all morphisms between them are necessarily finite; and consequently flat as well. To compute degrees, we may invert $N$ (because $\mathfrak{M}(S)$ is flat over $\Z$), and then the result is physically obvious \ref{cor:3.7.2}.
\end{proof}

\begin{kcor}\label{cor:5.5.4}
Let $E/S$ be an elliptic curve, $N > 1$ an integer. Then
\begin{enumerate}[(1)]
\item locally f.p.p.f. on $S$, any $\Gamma_1(N)$-structure can be completed to a balanced $\Gamma_1(N)$-structure, and any balanced $\Gamma_1(N)$-structure on $E/S$ can be completed to a $\Gamma(N)$-structure.
\item locally f.p.p.f. on $S$, $E/S$ admits $\Gamma_1(N)$, balanced $\Gamma_1(N)$, and $\Gamma(N)$-structures.
\item If $K$ is a finite locally free $S$-subgroup-scheme of rank $N$ in $E[N]$, and if $K$ is ``cyclic'' (i.e., locally f.p.p.f. admits a ``generator''), then the quotient $K'$ of $E[N]$ by $K$ is also cyclic.
\item a point $P \in E(S)$ has ``exact order $N$'' on $E/S$ if and only if locally f.p.p.f. on $S$ we can complete $P$ to a $\Gamma(N)$-structure $(P, Q)$ on $E/S$.
\end{enumerate}
\end{kcor}

\begin{proof}
Indeed (1) and (2) hold if we only allow finite locally free base-changes $S' \to S$; this is precisely the meaning of the previous corollary. To prove (3), we may f.p.p.f. localize on $S$ and suppose that $K$ is the subgroup-scheme ``generated'' by some point $P \in E(S)$ of ``exact order $N$.'' By (1), we may complete $P$ to balanced $\Gamma_1(N)$-structure $(P, K, Q')$, after a further finite locally free $S' \to S$. Having done so, we have produced a generator for $K'$, namely $Q'$. The final assertion (4) is just the concatenation of assertion (1) above and of assertion (1) of Theorem \ref{thm:5.5.2}.
\end{proof}

\begin{kcor}\label{cor:5.5.5}
Let $E/S$ be an elliptic curve, $N > 1$ an integer, $P \in E[N](S)$ a point of ``exact order $N$'' on $E/S$, $Q \in E[N](S)$ a point killed by $N$. Denote by $K \subset E[N]$ the cyclic subgroup-scheme of rank $N$ generated by $P$, by $E'$ the quotient of $E$ mod $K$, by $K' \subset E'$ the quotient of $E[N]$ mod $K$, and by $Q' \in K'(S)$ the image of $Q$. Then the following conditions are equivalent.
\begin{enumerate}[(1)]
\item $(P, Q)$ is a $\Gamma(N)$-structure on $E/S$
\item $(P, K, Q')$ is a balanced $\Gamma_1(N)$-structure on $E/S$.
\end{enumerate}
\end{kcor}

\begin{proof}
By (\ref{thm:5.5.2}), (1) implies (2). The converse is a special case of \ref{prop:1.11.3}, applied to the diagram
\[
\begin{tikzcd}[row sep=large, column sep=small]
0 \arrow[r] & \Z/N\Z \arrow[r] \arrow[d, "P"] & \Z/N\Z \oplus \Z/N\Z \arrow[r] \arrow[d, "{(P, Q)}"] & \Z/N\Z \arrow[r] \arrow[d, "Q'"] & 0 \\
0 \arrow[r] & K \arrow[r] & E[N] \arrow[r] & K' \arrow[r] & 0.
\end{tikzcd}
\]
\end{proof}

\begin{kcaut}\label{caut:5.5.6}
In the above corollary, consider the condition
\begin{enumerate}[(3)]
\item $Q'$ has ``exact order $N$'' in $E'/S$.
\end{enumerate}
Trivially (2) $\implies$ (3), but they are not equivalent. Indeed, if $N$ is a prime power $p^n$, and if $S$ is an $\F_p$-scheme, then the origin $Q' = 0$ has ``exact order $p^n$'' in $E'/S$, and $P = 0$ has ``exact order $p^n$'' on $E/S$. So if (3) were equivalent to (2), and so to (1), we would find that on any $E/S$ with $S/\F_p$, $(0,0)$ is a Drinfeld $p^n$-basis!
\end{kcaut}

\begin{kthm}\label{thm:5.5.7}
Let $E/S$ be an elliptic curve, $N > 1$ an integer, and $P, Q \in E[N](S)$ a pair of points of order $N$. Then
\begin{enumerate}[(1)]
\item If $(P, Q)$ is a Drinfeld $N$-basis of $E/S$, then for any divisor $d$ of $N$, the points $(dP, dQ)$ are a Drinfeld $N/d$-basis of $E/S$.
\item If $P$ has ``exact order $N$'' on $E/S$, then for every divisor $d$ of $N$, the point $dP$ has ``exact order $N/d$'' on $E/S$.
\item Suppose that $N = p^n$ with $p$ a prime, and $n \ge 2$. Then $(P, Q)$ is a Drinfeld $p^n$-basis of $E/S$ if and only if $(pP, pQ)$ is a Drinfeld $p^{n-1}$-basis on $E/S$.
\item If $N = p^n$ with $n \ge 2$ as in (3) above, then $P$ has ``exact order $p^n$'' on $E/S$ if and only if $pP$ has ``exact order $p^{n-1}$'' on $E/S$.
\end{enumerate}
\end{kthm}

\begin{proof}
The problem is local f.p.p.f. on $S$. By (4) of Corollary \ref{cor:5.5.4}, the assertions (2) and (4) are consequences of the assertions (1) and (3). The proof of (1) is entirely analogous to the proof of part (1) of Theorem \ref{thm:5.5.2}; it uses only the fact that $[\Gamma(N)]$ is flat over $\Z$, that $[\Gamma(N/d)]$ is relatively representable by a closed subscheme of the relatively affine moduli problem
\[
E/S \longmapsto E[N/d](S) \times E[N/d](S),
\]
and that the assertion is (by \ref{cor:3.7.2}) obviously true once we invert $N$.

To prove the ``if'' direction of (3), we must show that if $(pP, pQ)$ is a Drinfeld $p^{n-1}$-basis of $E/S$, then $(P, Q)$ is a Drinfeld $p^n$-basis of $E/S$. Let us temporarily denote by $\mathcal{P}$ the moduli problem
\[
\begin{array}{@{}l@{}}
\mathcal{P}(E/S) = \text{pairs of points } P, Q \text{ in } E[p^n](S) \\
\quad\text{such that } (pP, pQ) \text{ is a Drinfeld } p^{n-1}\text{-basis of } E/S.
\end{array}
\]
This is a relatively representable moduli problem; indeed for any $E/S$, we have a cartesian diagram of $S$-schemes
\[
\begin{tikzcd}[row sep=large, column sep=small]
(P, Q) \arrow[d] & \mathcal{P}_{E/S} \arrow[r, hook] \arrow[d] & E[p^n] \times E[p^n] \arrow[d, "{(p,p)}"', "{\shortstack{finite flat\\of degree $p^4$}}"] \\
(pP, pQ) & {[\Gamma(p^{n-1})]}_{E/S} \arrow[r, hook] & E[p^{n-1}] \times E[p^{n-1}]
\end{tikzcd}
\]
Therefore $\mathcal{P}$ is finite flat over $[\Gamma(p^{n-1})]$ of degree $p^4$, and hence it is finite and flat over $\Ell$, so in particular flat over $\Z$. By reduction to the universal case, it suffices to show that if $E/A$ is an elliptic curve over a ring $A$ which is flat over $\Z$, and if $P, Q$ in $E[p^n](A)$ are a $\mathcal{P}$-structure on $E/A$, then they are a $\Gamma(p^n)$-structure. The moduli problem $[\Gamma(p^n)]$ is relatively representable by a closed subscheme of the affine $A$-scheme $E[p^n] \times E[p^n]$:
\[
[\Gamma(p^n)]_{E/A} \hookrightarrow E[p^n] \times E[p^n].
\]
It is physically obvious (from \ref{cor:3.7.2}) that the pair $(P, Q)$ defines a $\Gamma(p^n)$-structure as soon as we pass to $A[1/p]$; therefore if $f_j = 0$ are the equations defining $[\Gamma(p^n)]$ inside $E[p^n] \times E[p^n]$, we have
\[
f_j(P, Q) = 0 \text{ in } A[1/p] \quad\text{whence}\quad f_j(P, Q) = 0 \text{ in } A
\]
because $A$ is flat over $\Z$.
\end{proof}

\begin{kcor}\label{cor:5.5.8}
Let $E/S$ be an elliptic curve, $N > 1$ an integer, $d$ a divisor of $N$ such that $N$ and $N/d$ have exactly the same prime divisors. Let $P, Q$ be a pair of points in $E[N](S)$. Then
\begin{enumerate}[(1)]
\item The pair $(P, Q)$ is a Drinfeld $N$-basis on $E/S$ if and only if the pair $(dP, dQ)$ is a Drinfeld $N/d$-basis on $E/S$.
\item The point $P$ has ``exact order $N$'' on $E/S$ if and only if the point $dP$ has ``exact order $N/d$'' on $E/S$.
\end{enumerate}
\end{kcor}

\begin{proof}
This is just the concatenation of (3) and (4) of Theorem \ref{thm:5.5.7}, with the elementary ``factorization into prime powers'' \ref{lem:3.5.1} for $\Gamma(N)$ and $\Gamma_1(N)$-structures.
\end{proof}

% ------------------------------------------------------------------------------
% Section 5.6
% ------------------------------------------------------------------------------
\section{Pairings}
\label{sec:5.6}

\begin{pitem}\label{lem:5.6.1}
In this section, we analyze the behavior of our moduli problems with respect to the pairing and the $\langle\ ,\ \rangle$ pairing (cf.\ \ref{sec:2.8}).
\end{pitem}

\begin{pitem}\label{lem:5.6.2}
Recall \ref{sec:1.12} that an $N$'th root of unity $\zeta \in \mu_N(S)$ is said to be a ``primitive $N$'th root of unity'' if it is a root of the $N$'th cyclotomic polynomial in the ring $\Gamma(S, \cO_S^\times)$, or equivalently, if viewed as a point of $\mathbb{G}_m(S)$ it has ``exact order $N$'' in $\mathbb{G}_m/S$, or equivalently if it ``generates'' $\mu_N(S)$.
\end{pitem}

\begin{kthm}\label{thm:5.6.3}
Let $E/S$ be an elliptic curve, $N > 1$ an integer. Then
\begin{enumerate}[(1)]
\item If $(P, Q)$ is a Drinfeld $N$-basis of $E/S$, then $e_N(P, Q)$ is a ``primitive $N$'th root of unity'' in $\mu_N(S)$.
\item If the data
\[
P \in (\Ker \pi)(S);\quad E \xrightarrow{\ \pi\ } E';\quad P' \in (\Ker \pi^t)(S)
\]
is a balanced $\Gamma_1(N)$-structure on $E/S$, then $\langle P, P'\rangle_\pi$ is a ``primitive $N$'th root of unity'' in $\mu_N(S)$.
\end{enumerate}
\end{kthm}

\begin{proof}
The question is f.p.p.f. local on $S$, so to prove (2) we may, by \ref{cor:5.5.4}, part (1), assume there is a Drinfeld $N$-basis on $E/S$, $(P, Q)$, with $P' = \pi Q$. Then $\langle P, P'\rangle_\pi = e_N(P, Q)$, so it suffices to prove (1).

Now (1) is physically obvious if $N$ is invertible. We will reduce to that case. By reduction to the universal case, we may assume that $S$ is affine, say $S = \Spec(A)$, with $A$ flat over $\Z$. We must show that
\[
\Phi_N(e_N(P, Q)) = 0 \text{ in } A,
\]
where $\Phi_N$ denotes the $N$'th cyclotomic polynomial. But we know that
\[
\Phi_N(e_N(P, Q)) = 0 \text{ in } A[1/N],
\]
and $A$ is flat over $\Z$.
\end{proof}

\begin{krem}\label{rem:5.6.4}
Given an elliptic curve $E/S$ with $N$ invertible on $S$, a pair of points $(P, Q)$ in $E[N](S)$ is a $\Gamma(N)$-structure if and only if $e_N(P, Q)$ is a primitive $N$'th root of unity in $\mu_N(S)$. This equivalence breaks down over a general $S$. Indeed if $N$ is a prime power $p^n$, then for any $\F_p$-scheme $S$, and any power $p^k$ of $p$, the trivial root of unity $1 \in \mu_{p^k}(S)$ is a ``primitive $p^k$'th root of unity.'' But in general $(0,0)$ is not a $\Gamma(p^k)$-structure on an elliptic curve over an $\F_p$-scheme, although $e_{p^k}(0,0) = 1$ is a ``primitive $p^k$'th root of unity.''
\end{krem}

\begin{kcor}\label{cor:5.6.5}
The moduli problems $[\Gamma(N)]$ and $[\mathrm{bal.}\ \Gamma_1(N)]$ are relatively representable over $\Ell$ by $\Z[\zeta_N]$-schemes, where $\Z[\zeta_N]$ denotes the ring $\Z[X]/(\Phi_N(X))$. Concretely, for any representable moduli problem $S$, the schemes
\[
\mathfrak{M}(S, \Gamma(N)),\quad \mathfrak{M}(S, \mathrm{bal.}\ \Gamma_1(N))
\]
have natural structures of schemes over $\mathfrak{M}(S)\otimes_\Z \Z[\zeta_N]$, defined by the functions $e_N(P, Q)$, $\langle P, P'\rangle_\pi$ on these schemes.
\end{kcor}

\begin{kterm}\label{term:5.6.6}
We will refer to the primitive $N$'th root of unity (resp.\ $\langle P, P'\rangle_\pi$) as the \emph{determinant} of the corresponding $\Gamma(N)$ (resp.\ balanced $\Gamma_1(N)$)-structure.
\end{kterm}
